%==============================================================================
% AAAI-27 submission — Austin-law benchmark
%
% BEFORE YOU COMPILE IN OVERLEAF:
%   1. Download the AAAI-27 Author Kit from aaai.org and drop `aaai2027.sty`
%      (and `aaai2027.bst`) into this folder. If the kit is still labelled by a
%      different year, rename the \usepackage line below to match the .sty you have
%      (e.g. aaai2026). The preamble below follows the kit's mandated form.
%   2. For anonymous review keep [submission]; remove it for camera-ready and fill
%      in the real \author block.
%   3. This file intentionally uses only amsmath/amsthm/booktabs/natbib so it drops
%      into Overleaf without extra dependencies. Cross-references are plain \ref.
%==============================================================================
\documentclass[letterpaper]{article}
\usepackage[submission]{aaai2027}  % <-- supply this .sty from the AAAI author kit

\usepackage{times}
\usepackage{helvet}
\usepackage{courier}
\usepackage[hyphens]{url}
\usepackage{graphicx}
\urlstyle{rm}
\def\UrlFont{\rm}
\usepackage{natbib}
\usepackage{caption}
\usepackage{amsmath}
\usepackage{amssymb}
\usepackage{amsthm}
\usepackage{booktabs}
\frenchspacing
\setlength{\pdfpagewidth}{8.5in}
\setlength{\pdfpageheight}{11in}

\pdfinfo{
/TemplateVersion (2027.1)
}

% --- theorem environments -----------------------------------------------------
\theoremstyle{definition}
\newtheorem{definition}{Definition}
\theoremstyle{plain}
\newtheorem{proposition}{Proposition}

% --- shorthands ---------------------------------------------------------------
\newcommand{\op}{\diamond}
\newcommand{\Law}{\mathrm{Law}}
\newcommand{\Eqtwo}{\mathrm{Eq}2}

\title{An Austin-Law Benchmark:\\ Verifiable, Renewable Construction Problems Beyond Automated Reach}

\author{
    %%%% anonymous for submission; fill in for camera-ready
    Anonymous Submission
}

\begin{document}
\maketitle

%==============================================================================
\begin{abstract}
The Equational Theories Project (ETP) is verifiable but uniform: a single strategy---
automated proof search together with a finite model builder---resolves twenty-two
million implications, so it measures execution rather than construction. We isolate the
sub-family on which that strategy is \emph{provably} inapplicable: \emph{Austin laws},
which admit infinite models but no nontrivial finite one and exist only at order five and
above, where finite model search is empty by theorem. We extend this family past the
order ETP enumerated and pose a two-sided, machine-checkable task: given a law certified
to have no nontrivial finite model, either construct a nontrivial model with a Lean proof
that it satisfies the law, or prove the law trivial. Exactly one side is achievable and
no instance is unanswerable. Against the strongest automated portfolio we can assemble
(Vampire, E, and Twee across term orderings and an order of magnitude of budget), the hard
tier yields \emph{zero} new nontrivial models, and the resolution rate is flat in the
compute budget---the frontier is bounded by method, not by time. The generator is public,
so the benchmark is minted rather than fixed. We report the corpus in equivalence classes
(a $\geq 26\%$ collapse), a Lean judge verified end to end, and two order-five open cases
from ETP confirmed to be Austin laws by two independent methods.
\end{abstract}

%==============================================================================
\section{Introduction}
\label{sec:intro}
% TODO — draft last, once the contribution list is frozen. Anchors:
% - ETP uniformity (one strategy resolves all) => measures execution not construction.
% - Austin laws: infinite models, no nontrivial finite; order >= 5; FMB empty by theorem.
% - Contributions: (i) two-sided task + certified answerability; (ii) order->=6 corpus +
%   public generator; (iii) Lean judge; (iv) strong baseline defining the hard tier;
%   (v) two order-5 open cases closed. Future work: the infinite-model formalisation.

%==============================================================================
\section{Background and Related Work}
\label{sec:related}
% TODO — write once claims are fixed so the debts are exact. Anchors:
% - ETP order<=4 enumerated; order-5 chapter (106 trivial-finite, 10 Austin, 96 open).
% - Austin (1983); Kisielewicz (28770, 5093).
% - JRS (saturation as explicit model): the honest debt. What it costs (existence is
%   mechanical is dead as a headline), what it gives (Twee, SZS/rewrite machinery).
% - Infinox (Claessen & Lilliestrom): our (i)-prover is a specialisation.
% - Berlioz & Mellies latent-space paper: adjacent, statistical, not a competitor.

%==============================================================================
\section{The Task}
\label{sec:task}

We work with magmas: a set $M$ equipped with a single binary operation
$\op : M \times M \to M$, with no further axioms. An \emph{equational law} is a
universally quantified identity $L : x = T[x, y, z, \dots]$, where $T$ is a term built
from the variables and $\op$. A magma \emph{satisfies} $L$ when the identity holds for
every assignment of its elements to the variables. Throughout, the fixed target of
interest is the law $\Eqtwo : x = y$, which a magma satisfies exactly when it is trivial
(has a single element); we write $L \models x = y$ for the semantic implication that
every model of $L$ is trivial.

The benchmark is built entirely from laws carrying a machine-checked certificate of a
single structural fact.

\begin{definition}[admissible instance]
\label{def:admissible}
A law $L$ is \emph{admissible} if there is a proof that $L$ has no nontrivial finite
model: every finite magma satisfying $L$ is trivial.
\end{definition}

This certificate is not a result of the benchmark; it is the ticket of admission, and it
is what makes the task well-posed. We produce it with a first-order argument---a subterm
of $T$ containing every occurrence of $x$ acts as a left inverse and is therefore
injective, so on a finite carrier it is surjective, and surjectivity of the relevant map
collapses the magma---mechanised as a query discharged by an automated prover. The
argument is a specialisation of Infinox's method for establishing finite
unsatisfiability~\citep{claessen2011infinox}, and we claim no novelty for it beyond the
observation that, for the syntactic shape $x = C[S(x)]$, the left inverse is available
without search.

Admissibility has a consequence that turns an open problem into an answerable question.

\begin{proposition}[dichotomy]
\label{prop:dichotomy}
For every admissible $L$, exactly one of the following holds: \emph{(a)}
$L \models x = y$, i.e.\ the implication into $\Eqtwo$ is valid and $L$ is
\emph{trivial}; or \emph{(b)} $L$ admits a nontrivial model, and every nontrivial model
of $L$ is infinite---$L$ is an \emph{Austin law}.
\end{proposition}

\begin{proof}
Immediate from admissibility together with the completeness of first-order logic: if
$L \not\models x = y$ then $L \cup \{\exists u\, v.\ u \neq v\}$ is consistent and so has
a model, which by Definition~\ref{def:admissible} cannot be finite.
\end{proof}

Austin laws---laws with infinite models but no nontrivial finite one---exist only at order
five and above~\citep{kisielewicz1988}, which is precisely why the finite regime cannot
exhibit them and why a finite model builder, one of the two standard tools, is
inapplicable to this family by theorem rather than merely slow.

The dichotomy licenses a two-sided task. \textbf{Given an admissible law $L$, either
exhibit a nontrivial model of $L$ together with a machine-checkable proof that it
satisfies $L$, or prove $L \models x = y$.} Exactly one side is achievable, both are
machine-checkable, and---this is the point that distinguishes the construction here from
an open-problems list---\emph{no instance is unanswerable}. The admissibility certificate,
combined with Proposition~\ref{prop:dichotomy}, guarantees that every instance has a
determinate answer of one of the two kinds. This is a stronger guarantee than a benchmark
can usually offer: the difficulty of an instance is bounded below by the effort a solver
must expend, never by the possibility that the instance is ill-posed or the answer
nonexistent.

We stress what is deliberately \emph{not} a design criterion: solvability. An instance
that no current method resolves still discriminates between methods, and an instance whose
answer is unreachable in principle would be a defect, but the two are far apart. Every
answer here is a finite object---a Lean proof---and the space of Lean proofs is
recursively enumerable, so no admissible instance is unsolvable in principle; it is only
hard relative to a fixed, named method. Hardness relative to a published baseline is the
normal and sufficient notion for a benchmark, and it is the one we adopt
(the Baseline section). Conversely, the benchmark is not vacuous merely because its hard
tier is unresolved: the certificate guarantees that the resolution exists.

%==============================================================================
\section{Verification}
\label{sec:verification}

A benchmark of construction problems is only as meaningful as its ability to check a
proposed construction. We make the checker unconditional by fixing a single arbiter: an
answer is a proof in the Lean proof assistant of a statement we generate mechanically from
the law, and the proof is accepted exactly when Lean's kernel accepts it and its logical
dependencies lie within a fixed allowlist. Nothing about the \emph{form} of the model---
whether it is presented as a subset of the integers, a quotient of a term algebra, or a
rewrite system---enters the judge; any Lean proof of the generated goal is a valid answer.

For a law $L : x = T$ with variables $x, y, z, \dots$, the judge emits the definitions
\begin{align}
\Law(\op) &\equiv \forall x\, y\, z\, \dots \in M,\ x = T[\op], \label{eq:law}\\
\textsc{AustinGoal} &\equiv \exists\, (M : \mathrm{Type})\, (\op : M \to M \to M), \notag\\
&\qquad (\exists\, a\, b : M,\ a \neq b) \wedge \Law(\op), \label{eq:austingoal}\\
\textsc{TrivialGoal} &\equiv \forall\, (M : \mathrm{Type})\, (\op : M \to M \to M), \notag\\
&\qquad \Law(\op) \to \forall\, a\, b : M,\ a = b. \label{eq:trivialgoal}
\end{align}
A submission proves one of the two goals. The nontriviality clause
$\exists a\, b.\ a \neq b$ is what forces an Austin answer to exhibit two distinct
elements---and, together with admissibility, what forces the witnessing magma to be
infinite---while keeping the statement itself finiteness-free, so that a submitted model
is a genuine model regardless of whether the admissibility prover was correct.

The judged artifact is the concatenation of a generated header, the submitter's body, and
a generated footer. The header fixes the statement; the footer closes with
\texttt{example : Problem.AustinGoal := solution}, which forces the submitter's
\texttt{solution} to elaborate at exactly the type the header defines. The submitter never
writes the statement and therefore cannot weaken it. Two safeguards make this airtight.
First, a textual gate run before Lean rejects the escape hatches that would let a file
typecheck without proving the intended proposition: \texttt{sorry} and \texttt{admit};
\texttt{native\_decide}, which discharges goals by trusted compiled code outside the
kernel; user-declared \texttt{axiom}s; the metaprogramming constructs (\texttt{macro},
\texttt{syntax}, \texttt{elab}) that could redefine notation; and, most importantly, any
redefinition of the generated names, since the characteristic attack is to shadow
\textsc{AustinGoal} with a trivially true proposition and prove that instead. Second, after
Lean accepts the file, the judge reads the axiom footprint printed by
\texttt{\#print axioms solution} and requires it to be a subset of
$\{\texttt{propext}, \texttt{Quot.sound}, \texttt{Classical.choice}\}$; the appearance of
\texttt{sorryAx} or the \texttt{native\_decide} axiom, or of any axiom outside the
allowlist, is a rejection even when the file compiled. The allowlist is closed rather than
a blocklist of known cheats, so an unforeseen axiom fails by default.

We have verified the judge end to end against a reference answer---a two-element magma
satisfying a simple admissible law---which passes the textual gate, elaborates against the
generated goal, and exhibits the expected axiom footprint; the negative controls (a
\texttt{sorry}ed proof, an axiom smuggle, and both the namespaced and bare forms of the
shadowing attack) are each rejected before Lean runs.

Two channels for producing accepted proofs already exist. Concrete algebraic models---a
finite ring, a $\mathbb{Z}$-module, an adjoined-root extension---are checked directly: the
model's operation is a Lean definition and the law is discharged either by a symbolic
identity or, for small carriers, by kernel computation. This is the channel by which a
previously open order-five instance was closed by a mod-$17$ affine model
(the Closed Cases section). The second channel, and the one that reaches the hardest
instances, presents the model as the term algebra modulo a convergent rewrite system
derived from a prover saturation; certifying such a model in Lean requires formalising
ground confluence of ordered rewriting, which we have reduced to a single open lemma in an
otherwise complete development and which we treat as the subject of a companion paper. The
judge itself is agnostic to this: it is sound today for every answer a submitter can drive
through the kernel, and widening the set of \emph{easily producible} answers is a matter
of building solver-side tools, not of changing the arbiter.

%==============================================================================
\section{The Corpus and the Generator}
\label{sec:corpus}
% TODO — draft. Anchors: one-op extension (~500x yield); >=26% collapse (262 -> <=195
% classes) by prover-free model separation (1123/1128); quote classes; seed_dedupe;
% contamination-free (order<=4 public, order-5 table, ship order>=6 with seeds+dates);
% durability = classes per thousand surviving extensions.

%==============================================================================
\section{The Baseline and the Hard Tier}
\label{sec:baseline}

The hard tier of the benchmark is not a fixed list of laws but a property relative to a
named method: an admissible law belongs to the hard tier when a published, reproducible
portfolio of automated provers fails to resolve it---fails to prove it trivial and fails
to construct a nontrivial model---within a stated budget. This section defines that
portfolio, reports how the tier behaves under it, and establishes the two facts the
benchmark rests on: that the frontier is bounded by \emph{method} rather than by
\emph{compute}, and that it is not an artifact of any single prover.

\subsection{The portfolio}
A single prover under a single term ordering is not a baseline; it is a shrug. A law whose
completion diverges under one reduction ordering may terminate under another, so a
defensible baseline must range over provers, orderings, and budget. Portfolio v1.0
comprises eight configurations: Vampire~5.0.1 in five modes---three proof-search
configurations targeting $L \models x = y$ and two saturation configurations testing
consistency, under both KBO and LPO orderings---the E prover in a proof and a saturation
mode, and Twee~2.6.1, an implementation of unfailing completion purpose-built for the unit
equational fragment in which every law of our corpus lives. Each configuration is run over
a budget ladder of $30, 60, 120, 300,$ and $600$ seconds, and a law is recorded as
resolved as soon as any configuration returns a verdict, whereupon the remaining ladder is
skipped.

The two directions of the task are found by different machinery, and the portfolio must
contain both: a proof of triviality is proof search, whereas a nontrivial model is
witnessed by a saturation closing---a saturated clause set in this fragment \emph{is} a
ground-convergent rewrite system and hence an explicit model~\citep{janota2026saturations}.
A portfolio that only refutes would never find the second kind of answer. Every verdict is
mapped from the prover's SZS status line, and a \texttt{Satisfiable} verdict from a
strategy that sacrificed completeness is discarded rather than trusted. Provers, versions,
and exact flags are pinned in a container image released with the benchmark, and the
harness refuses to run unless a known-answer self-test passes.

\subsection{The tier is method-bound, not compute-bound}
The central empirical question is whether the hard tier is genuinely beyond the reach of
these methods or merely under-budgeted. The two are separable by measurement: run the
ladder and record, at each budget, the fraction of laws resolved. If resolution keeps
rising with budget, the frontier is a matter of compute; if it flattens, the surviving
laws are bound by what the methods can express, not by how long they run.

Over a sample of 120 laws from the no-finite-model tier, the portfolio resolves three laws
at 30 seconds and not one additional law at any higher budget (Table~\ref{tab:curve}).
Twenty times the compute buys nothing. This licenses the definition we adopt: the hard
tier is the set of admissible laws left unresolved at the ladder's maximum budget under
\emph{every} configuration, and the flatness of the curve past 30 seconds is the evidence
that this maximum lies beyond the knee rather than short of it. We report the curve, not a
single timeout, precisely so this judgment is checkable.

\begin{table}[t]
\centering
\begin{tabular}{lcc}
\toprule
budget & laws resolved & $\Delta$ rate \\
\midrule
$30$\,s  & $3/120 = 0.025$ & --- \\
$60$\,s  & $0/117 = 0.000$ & $-0.025$ \\
$120$\,s & $0/117 = 0.000$ & $0$ \\
$300$\,s & $0/117 = 0.000$ & $0$ \\
$600$\,s & $0/117 = 0.000$ & $0$ \\
\bottomrule
\end{tabular}
\caption{Resolution rate of portfolio v1.0 over the budget ladder, on a 120-law sample of
the no-finite-model tier. Every law that resolves, resolves at 30\,s; higher budgets add
nothing. Zero laws acquire a nontrivial model at any budget.}
\label{tab:curve}
\end{table}

A prior single-prover measurement points the same way and rules out the most deflationary
explanation. Re-running the corpus's earlier classifier at 300 seconds per prover, an
order of magnitude above its original budget, converted $3.7\%$ of laws---and of the 216
no-finite-model laws in that pass, four became provably trivial while \emph{none} acquired
a nontrivial model. The laws that resisted the shorter budget resisted the longer one on
the construction side without exception; the only movement was the removal of trivial laws
the shorter budget had failed to prove trivial.

\subsection{The frontier is not an artifact of one prover}
Flatness under a portfolio could in principle hide a weakness common to all its members.
The sharpest test available is Twee, whose completion procedure is algorithmically
distinct from Vampire's and E's superposition and is the tool most likely, on prior
grounds, to complete a system the others leave open. It does complete systems the others
find awkward---on one of the two order-five laws we later close, Vampire's saturation runs
to 357 clauses while Twee's unfailing completion terminates in seconds. Yet across the
sampled tier, Twee resolves \emph{no} law that the rest of the portfolio does not, and in
particular contributes no new nontrivial model. The reshaping outcome we explicitly tested
for---a law solved only by completion, which would have shown the tier to be an artifact
of Vampire's search rather than a genuine frontier---does not occur. Twee is a strong
member of the portfolio, not a key to the tier.

\subsection{Contamination and what resolution means here}
\label{sec:contamination}
Every law the portfolio does resolve, it resolves as \emph{trivial}: of the sampled tier,
eleven laws are proved to satisfy $L \models x = y$, and zero acquire a nontrivial model.
These trivial laws are contamination rather than error. They are admissible---they have no
nontrivial finite model---and the shorter classifier budget had failed to prove them
trivial, so they sat in the no-finite-model tier without being Austin. The stronger
portfolio removes them, and the two-sided task is exactly what makes this harmless: a law
that turns out trivial is a legitimately answered instance, not a mislabelled one. The
completion provers carry more of this load than Vampire's proof modes---one law is proved
trivial by E in one second and by Twee in twelve while all five Vampire configurations
exhaust thirty seconds---which is a second, quieter argument for a portfolio over a single
strong prover.

That the portfolio produces zero new models is therefore the result to hold onto. The hard
tier's construction side is untouched by the strongest automated methods we can assemble,
across orderings and an order of magnitude of budget, while every instance in it is
certified to possess an answer (the Task section). The gap between \emph{an answer exists}
and \emph{this portfolio finds it} is the quantity the benchmark measures.

\subsection{Scope of the measurement}
The 120-law figure is a sample of the roughly 3{,}400 no-finite-model laws in the current
corpus, and it is sound for a \emph{rate}: the shape of the resolution curve and the
absence of completion-only models are properties of the population, estimated without bias
by the sample. It is not the tier's \emph{membership list}. Publishing the definitive hard
tier---the exact laws that survive portfolio v1.0 at maximum budget---requires the full
sweep, a single long run we schedule once the ladder and portfolio are frozen, since its
cost is dominated by the laws that never resolve and therefore traverse the entire ladder.
The membership list refines the artifact; it does not change the method or the rate
reported here.

%==============================================================================
\section{Two Closed Order-Five Cases}
\label{sec:closed-cases}
% TODO — draft. Anchors: 12857, 33436 in ETP Table 20.2 (open); confirmed Austin two ways
% (Vampire saturation + Twee, both CounterSatisfiable); models computable & non-vacuous
% (27/27, rules fire, x=y refuted); neither is a plain TRS even after Twee reduction
% => CSI/TTT2 inapplicable, Lean cert needs ground confluence (companion paper).

%==============================================================================
\section{Discussion and Limitations}
\label{sec:discussion}
% TODO. Anchors: hard-tier collapse not yet measured (prover-only census there);
% ground_confluent open => saturation-model answers not yet Lean-certifiable (algebraic
% ones are), framed as next paper; optional hand-solved demonstrations; (i)-prover vs
% Infinox not yet benchmarked.

%==============================================================================
\section{Conclusion}
\label{sec:conclusion}
% TODO. A public generator + certified two-sided task => the benchmark is minted, not
% fixed, and hard relative to a named, published baseline.

%==============================================================================
\bibliographystyle{aaai2027}  % supply aaai2027.bst from the author kit
\bibliography{references}

\end{document}
